\section*{Chapitre 4 \textemdash\ \'Equations et In\'equations}
\addcontentsline{toc}{section}{Chapitre 4 \textemdash\ \'Equations et In\'equations}

\subsection*{Exercices}
\addcontentsline{toc}{subsection}{Exercices}

%────────────────────────────────────────────────────────────────
\solutiontitle{4}{1}{R\'esoudre avec le discriminant}

\begin{enumerate}
  \item $x^2-5x+6=0$~: $\Delta=25-24=1$~; $x=\frac{5\pm1}{2}$~: $x_1=3$, $x_2=2$.
  \item $2x^2-3x+1=0$~: $\Delta=9-8=1$~; $x=\frac{3\pm1}{4}$~: $x_1=1$, $x_2=\frac{1}{2}$.
  \item $x^2+2x+5=0$~: $\Delta=4-20=-16<0$. Pas de racine r\'eelle.
  \item $4x^2-12x+9=0$~: $\Delta=144-144=0$. Racine double~: $x_0=\frac{3}{2}$.
  \item $x^2-7=0$~: $x=\pm\sqrt{7}$.
  \item $3x^2+x=0$~: $x(3x+1)=0$~; $x=0$ ou $x=-\frac{1}{3}$.
  \item $x^2-2x-8=0$~: $\Delta=4+32=36$~; $x=\frac{2\pm6}{2}$~: $x_1=4$, $x_2=-2$.
  \item $5x^2+3x-2=0$~: $\Delta=9+40=49$~; $x=\frac{-3\pm7}{10}$~: $x_1=\frac{2}{5}$, $x_2=-1$.
\end{enumerate}

%────────────────────────────────────────────────────────────────
\solutiontitle{4}{2}{Forme canonique}

\begin{enumerate}
  \item $x^2-6x+11=(x-3)^2+2$. Sommet $(3,2)$.
  \item $2x^2+4x-1=2(x+1)^2-3$. Sommet $(-1,-3)$.
  \item $-x^2+2x+3=-(x-1)^2+4$. Sommet $(1,4)$.
  \item $x^2+x+1=\bigl(x+\tfrac{1}{2}\bigr)^2+\tfrac{3}{4}$.
  \item $3x^2-6x+5=3(x-1)^2+2$.
  \item $-2x^2+8x-7=-2(x-2)^2+1$.
\end{enumerate}

%────────────────────────────────────────────────────────────────
\solutiontitle{4}{3}{In\'equations du second degr\'e}

\begin{enumerate}
  \item $x^2-x-2\leq0$~: racines $-1$ et $2$ ($a>0$). Solution~: $[-1,2]$.
  \item $x^2+1>0$~: toujours vraie. Solution~: $\R$.
  \item $-x^2+4x-3\geq0$~: racines $1$ et $3$ ($a=-1<0$), positif entre elles.
        Solution~: $[1,3]$.
  \item $2x^2-5x+3<0$~: $\Delta=1$, racines $1$ et $\frac{3}{2}$ ($a>0$).
        Solution~: $\bigl]1,\tfrac{3}{2}\bigr[$.
  \item $x^2-4\geq0$~: racines $\pm2$. Solution~: $]-\infty,-2]\cup[2,+\infty[$.
  \item $x^2-4x+4=(x-2)^2\leq0$~: uniquement en $x=2$. Solution~: $\{2\}$.
\end{enumerate}

%────────────────────────────────────────────────────────────────
\solutiontitle{4}{4}{Relations de Vi\`ete}

\begin{enumerate}
  \item $x^2-7x+10$~: somme $S=7$, produit $P=10$.
  \item $3x^2+x-2$~: $S=-\tfrac{1}{3}$, $P=-\tfrac{2}{3}$.
  \item $-x^2+4x-4$ ($a=-1$)~: $S=-b/a=4$, $P=c/a=4$.
  \item $x^2-5$~: $S=0$, $P=-5$.
  \item $2x^2-6x$~: $S=3$, $P=0$.
  \item $x^2-2x+1$~: $S=2$, $P=1$.
\end{enumerate}

%────────────────────────────────────────────────────────────────
\solutiontitle{4}{5}{V\'erifier qu'un nombre est racine}

\begin{enumerate}
  \item $P(1)=1-6+11-6=0$.\checkmark
  \item $Q(1)=1+2-1-2=0$.\checkmark
  \item $R(2)=16-4-10-2=0$.\checkmark
  \item $S(-2)=-8-3(-2)+2=-8+6+2=0$.\checkmark
\end{enumerate}

%────────────────────────────────────────────────────────────────
\solutiontitle{4}{6}{Factoriser un polyn\^ome de degr\'e 3}

\textbf{(1)} $P(1)=1-6+11-6=0$.\checkmark\quad
$P(x)=(x-1)(x^2-5x+6)=(x-1)(x-2)(x-3)$.

\textbf{(2)} $Q(1)=0$.\checkmark\quad
$Q(x)=(x-1)(x^2+3x+2)=(x-1)(x+1)(x+2)$.

\textbf{(3)} $R(3)=27-27-12+12=0$.\checkmark\quad
$(x-3)(x^2+ex+f)$~: $e=0$, $f=-4$~; $R(x)=(x-3)(x-2)(x+2)$.

\textbf{(4)} $S(2)=16-4-10-2=0$.\checkmark\quad
$S(x)=(x-2)(2x^2+3x+1)=(x-2)(2x+1)(x+1)$.

%────────────────────────────────────────────────────────────────
\solutiontitle{4}{7}{Ensemble de d\'efinition d'une fraction rationnelle}

\begin{enumerate}
  \item $D=\R\setminus\{3\}$.
  \item $x^2-1=(x-1)(x+1)$~: $D=\R\setminus\{-1,1\}$.
  \item $x^2+1>0$ toujours~: $D=\R$.
  \item $x^2-4=(x-2)(x+2)$~: $D=\R\setminus\{-2,2\}$.
  \item $(x-1)(x+2)=0\Leftrightarrow x=1$ ou $x=-2$~: $D=\R\setminus\{-2,1\}$.
  \item $x^2-3x+2=(x-1)(x-2)$~: $D=\R\setminus\{1,2\}$.
\end{enumerate}

%────────────────────────────────────────────────────────────────
\solutiontitle{4}{8}{Simplifier une fraction rationnelle}

\begin{enumerate}
  \item $\dfrac{x^2-4}{x-2}=x+2$ pour $x\ne2$.
  \item $\dfrac{x^2-1}{x+1}=x-1$ pour $x\ne-1$.
  \item $\dfrac{x^2-3x+2}{x-1}=x-2$ pour $x\ne1$.
  \item $\dfrac{x^2+2x}{x}=x+2$ pour $x\ne0$.
  \item $\dfrac{(x-1)^2}{x^2-1}=\dfrac{x-1}{x+1}$ pour $x\ne\pm1$.
  \item $\dfrac{x^3-x}{x^2-1}=x$ pour $x\ne\pm1$.
\end{enumerate}

%────────────────────────────────────────────────────────────────
\solutiontitle{4}{9}{R\'esoudre un syst\`eme $2\times2$}

\begin{enumerate}
  \item $y=5-2x$~; $x-3(5-2x)=-1\Rightarrow7x=14$~: $x=2$, $y=1$.
  \item $y=3x-1$~; $x+2(3x-1)=12\Rightarrow7x=14$~: $x=2$, $y=5$.
  \item Somme~: $3x=9\Rightarrow x=3$~; $y=4$.
  \item $y=10-2x$~; $4x-3(10-2x)=0\Rightarrow10x=30$~: $x=3$, $y=4$.
\end{enumerate}

%────────────────────────────────────────────────────────────────
\solutiontitle{4}{10}{Tableau de signe d'un trin\^ome}

\begin{enumerate}
  \item $x^2-4x+3>0$~: racines $1$ et $3$ ($a>0$). Solution~: $]-\infty,1[\cup]3,+\infty[$.
  \item $-2x^2+x+1\geq0$~: $\Delta=9$, racines $-\frac{1}{2}$ et $1$ ($a<0$).
        Solution~: $[-\frac{1}{2},1]$.
  \item $x^2-2x+1=(x-1)^2\leq0$~: Solution~: $\{1\}$.
  \item $(2x-1)(x+3)<0$~: z\'eros $-3$ et $\frac{1}{2}$.
        Solution~: $]-3,\frac{1}{2}[$.
\end{enumerate}

%────────────────────────────────────────────────────────────────
\solutiontitle{4}{11}{Construire une \'equation de solution donn\'ee}

On utilise $x^2-(x_1+x_2)x+x_1x_2=0$.
\begin{enumerate}
  \item $x^2-7x+10=0$.
  \item $x^2-2x-3=0$.
  \item $(x-4)^2=0$, soit $x^2-8x+16=0$.
  \item $S=0$, $P=-2$~: $x^2-2=0$.
  \item $S=2$, $P=1-3=-2$~: $x^2-2x-2=0$.
  \item $x^2-7x=0$.
\end{enumerate}

%────────────────────────────────────────────────────────────────
\solutiontitle{4}{12}{Valeur num\'erique d'un polyn\^ome}

\textbf{(1)} $P(x)=x^3-3x^2+2x-1$~:
$P(-1)=-7$~; $P(0)=-1$~; $P(2)=-1$~; $P(3)=5$.

\textbf{(2)} $P(x)=2x^3+x^2-x+4$~:
$P(-1)=4$~; $P(0)=4$~; $P(2)=22$~; $P(3)=64$.

%────────────────────────────────────────────────────────────────
\solutiontitle{4}{13}{Discriminant et nombre de solutions}

\begin{enumerate}
  \item $\Delta=25-36=-11<0$~: aucune racine r\'eelle.
  \item $\Delta=36-36=0$~: une racine double.
  \item $\Delta=49+32=81>0$~: deux racines distinctes.
  \item $\Delta=1-4=-3<0$~: aucune racine r\'eelle.
  \item $\Delta=16-16=0$~: une racine double.
  \item $\Delta=4+60=64>0$~: deux racines distinctes.
\end{enumerate}

%────────────────────────────────────────────────────────────────
\solutiontitle{4}{14}{Lire les racines sur un graphe}

\begin{enumerate}
  \item $f(x)=x^2-4x+3$~: $\Delta=4$, racines $x=1$ et $x=3$.
        Sommet $(2,-1)$~; parabole ouverte vers le haut.
  \item $g(x)=-x^2+2x+3$~: $\Delta=16$, racines $x=-1$ et $x=3$.
        Sommet $(1,4)$~; parabole ouverte vers le bas.
  \item V\'erification alg\'ebrique~: les deux calculs de $\Delta$ confirment
        les lectures graphiques.\hfill$\square$
\end{enumerate}

%────────────────────────────────────────────────────────────────
\solutiontitle{4}{15}{Signe d'une fraction simple}

\begin{enumerate}
  \item $\frac{x-2}{x+1}$~: signe $+$ sur $]-\infty,-1[$, $-$ sur $]-1,2[$,
        $+$ sur $]2,+\infty[$.
  \item $\frac{-x+3}{x-5}$~: num. nul en $3$, d\'enom. nul en $5$.
        Signe~: $-$, $+$, $-$ sur les trois intervalles.
  \item $\frac{x^2-1}{x}=\frac{(x-1)(x+1)}{x}$~: z\'eros $-1$, $0$ (exclu), $1$.
        Signe~: $-,+,-,+$ sur les quatre intervalles.
  \item $\frac{(x+2)^2}{x-4}$~: $(x+2)^2\geq0$. Signe de $(x-4)$~:
        n\'egatif sur $]-\infty,4[$, positif sur $]4,+\infty[$, nul en $x=-2$.
\end{enumerate}

%────────────────────────────────────────────────────────────────
\solutiontitle{4}{16}{Fraction rationnelle~: signe et domaine}

\begin{enumerate}
  \item $F(x)=\frac{(x-2)(x+1)}{(x-2)(x+2)}=\frac{x+1}{x+2}$ ($x\ne2,-2$).
        Z\'ero en $-1$, exclu en $-2$ et $2$.
        $F>0$~: $]-\infty,-2[\cup]-1,2[\cup]2,+\infty[$.

  \item $F(x)=\frac{(x-1)(x+2)}{(x-3)(x+3)}$. $D=\R\setminus\{-3,3\}$, z\'eros~: $1$ et $-2$.
        $F>0$~: $]-\infty,-3[\cup]-2,1[\cup]3,+\infty[$.

  \item $F(x)=\frac{x(x-3)}{(x-3)(x+2)}=\frac{x}{x+2}$ ($x\ne3,-2$).
        Z\'ero en $0$.
        $F>0$~: $]-\infty,-2[\cup]0,3[\cup]3,+\infty[$.

  \item $F(x)=\frac{x(x-1)(x+1)}{x^2+1}$~: $D=\R$ (d\'enom. $>0$).
        Z\'eros~: $-1,0,1$. $F>0$~: $]-1,0[\cup]1,+\infty[$.
\end{enumerate}

%────────────────────────────────────────────────────────────────
\solutiontitle{4}{17}{Factoriser un polyn\^ome de degr\'e 3 compl\`etement}

On cherche une racine enti\`ere par test.

\textbf{(1)} $P(1)=1-7+6=0$. $P(x)=(x-1)(x^2+x-6)=(x-1)(x-2)(x+3)$.

\textbf{(2)} $Q(2)=8+4-8-4=0$. $Q(x)=(x-2)(x^2+3x+2)=(x-2)(x+1)(x+2)$.

\textbf{(3)} $R(3)=54-27-33+6=0$. $R(x)=(x-3)(2x^2+3x-2)=(x-3)(2x-1)(x+2)$.

\textbf{(4)} $S(-1)=-1-4-1+6=0$. $S(x)=(x+1)(x^2-5x+6)=(x+1)(x-2)(x-3)$.

%────────────────────────────────────────────────────────────────
\solutiontitle{4}{18}{Tableau de signe d'un polyn\^ome de degr\'e 3}

\begin{enumerate}
  \item $(x-1)(x-2)(x-3)\geq0$~: racines $1,2,3$.
        Solution~: $[1,2]\cup[3,+\infty[$.
  \item $(x-1)(x^2+x+1)>0$~: $\Delta(x^2+x+1)=-3<0$, donc $x^2+x+1>0$ toujours.
        Signe = signe de $(x-1)$. Solution~: $]1,+\infty[$.
  \item $x^3+x^2-4x-4=(x-2)(x+1)(x+2)\leq0$ (voir exo~17).
        Racines~: $-2,-1,2$. Solution~: $]-\infty,-2]\cup[-1,2]$.
  \item $2x^3-x^2-x=x(2x+1)(x-1)\geq0$~: racines $-\frac{1}{2},0,1$.
        Solution~: $[-\frac{1}{2},0]\cup[1,+\infty[$.
\end{enumerate}

%────────────────────────────────────────────────────────────────
\solutiontitle{4}{19}{Op\'erations sur les fractions rationnelles}

\begin{enumerate}
  \item $\frac{x^2-1}{x+1}+\frac{1}{x-1}=(x-1)+\frac{1}{x-1}
        =\frac{(x-1)^2+1}{x-1}=\frac{x^2-2x+2}{x-1}$ \quad($x\ne\pm1$).
  \item $\frac{x}{x-2}-\frac{2}{x+1}
        =\frac{x(x+1)-2(x-2)}{(x-2)(x+1)}=\frac{x^2-x+4}{(x-2)(x+1)}$.
  \item $\frac{x^2-4}{x-2}\times\frac{x-1}{x+2}=(x+2)\cdot\frac{x-1}{x+2}=x-1$\quad($x\ne\pm2$).
  \item $\frac{x^2-9}{x+3}\div(x-3)=\frac{(x-3)(x+3)}{x+3}\cdot\frac{1}{x-3}=1$\quad($x\ne\pm3$).
  \item $\frac{1}{x}+\frac{1}{x+1}+\frac{1}{x+2}
        =\frac{(x+1)(x+2)+x(x+2)+x(x+1)}{x(x+1)(x+2)}=\frac{3x^2+6x+2}{x(x+1)(x+2)}$.
  \item $\left(\frac{x}{x+1}\right)^2-\frac{x^2}{x^2-1}
        =\frac{x^2}{x+1}\!\left(\frac{1}{x+1}-\frac{1}{x-1}\right)
        =\frac{-2x^2}{(x+1)^2(x-1)}$\quad($x\ne\pm1$).
\end{enumerate}

%────────────────────────────────────────────────────────────────
\solutiontitle{4}{20}{Param\`etre dans le discriminant}

\textbf{(1)} $\Delta=m^2-4m=m(m-4)$.
Deux racines~: $m<0$ ou $m>4$.
Racine double~: $m=0$ ou $m=4$.
Pas de racine~: $0<m<4$.

\textbf{(2)} $\Delta=4(1-m)>0\Leftrightarrow m<1$ (et $m\ne0$).

\textbf{(3)} $1-(m+1)+m=0$~: toujours vrai. Donc $1$ est toujours racine.
Par Vi\`ete~: $x_2=m$.

\textbf{(4)} Conditions cumulatives~:
$x_1+x_2=-2m>0\Rightarrow m<0$~;
$x_1x_2=m+2>0\Rightarrow m>-2$~;
$\Delta=4(m-2)(m+1)\geq0\Rightarrow m\leq-1$ ou $m\geq2$.
Intersection~: $\mathbf{-2<m\leq-1}$.\hfill$\square$

%────────────────────────────────────────────────────────────────
\solutiontitle{4}{21}{Syst\`eme m\^el\'e lin\'eaire/quadratique}

\textbf{(1)} $x^2-2x=x+4\Rightarrow x^2-3x-4=(x-4)(x+1)=0$.
Solutions~: $(4,8)$ et $(-1,3)$.

\textbf{(2)} $x$ et $y$ sont racines de $t^2-5t+6=(t-2)(t-3)=0$.
Solutions~: $(2,3)$ et $(3,2)$.

\textbf{(3)} $x^2+(x+1)^2=25\Rightarrow x^2+x-12=0$~; $\Delta=49$.
$x=3,y=4$ ou $x=-4,y=-3$.

%────────────────────────────────────────────────────────────────
\solutiontitle{4}{22}{R\'esoudre une \'equation avec fraction rationnelle}

\begin{enumerate}
  \item $x\ne1$~: $x+1=2(x-1)\Rightarrow x=3$.
  \item $x\ne0,-2$~: $x^2+(x+2)=x(x+2)\Rightarrow x=2$.
  \item $x\ne-1$~: $x-1=x-2\Rightarrow-1=-2$. \textbf{Pas de solution.}
  \item $x\ne0,2$~: $3x=(x-2)(x+1)\Rightarrow x^2-4x-2=0$~; $x=2\pm\sqrt{6}$.
\end{enumerate}

%────────────────────────────────────────────────────────────────
\solutiontitle{4}{23}{In\'egalit\'e avec fraction rationnelle}

\begin{enumerate}
  \item $\frac{5}{x-2}>0\Rightarrow x>2$. Solution~: $]2,+\infty[$.
  \item $\frac{(x-2)(x+2)}{x+1}\leq0$~: z\'eros $\pm2$, exclu $-1$.
        Solution~: $]-\infty,-2]\cup]-1,2]$.
  \item $\frac{-4x}{x^2-1}\geq0$~: z\'eros $-1$ (exclu), $0$, $1$ (exclu).
        Solution~: $]-\infty,-1[\cup[0,1[$.
  \item $\frac{x+1}{x}>0$~: z\'ero $-1$, exclu $0$.
        Solution~: $]-\infty,-1[\cup]0,+\infty[$.
\end{enumerate}

%────────────────────────────────────────────────────────────────
\solutiontitle{4}{24}{Application~: probl\`eme de g\'eom\'etrie}

\textbf{(1)} $l+L=10$, $lL=24$~: $\Delta=100-96=4$~; dimensions $4\times6$~cm.

\textbf{(2)} $l(l+5)=204\Rightarrow l^2+5l-204=0$~; $\Delta=841=29^2$~;
$l=12$~m, $L=17$~m.

\textbf{(3)} $a^2=(a-2)(a+3)\Rightarrow a=6$.
Carr\'e $6\times6$, rectangle $4\times9$ (m\^eme aire $36$).\checkmark

%────────────────────────────────────────────────────────────────
\solutiontitle{4}{25}{Division euclidienne de polyn\^omes}

\textbf{(1)} $P(1)=1-2+1-1=-1$. Division~: $x^3-2x^2+x-1=(x-1)(x^2-x)-1$.

\textbf{(2)} $P(-2)=2(-8)+4+6+2=0$. Division~: $(x+2)(2x^2-3x+1)=(x+2)(2x-1)(x-1)$.

\textbf{(3)} $P(1)=0$. Division~: $(x-1)(x^2+x+2)$.
V\'erif.~: $\Delta(x^2+x+2)=-7<0$~: pas de factorisation r\'eelle suppl\'ementaire.\hfill$\square$

%────────────────────────────────────────────────────────────────
\solutiontitle{4}{26}{Factoriser ou d\'evelopper}

\begin{enumerate}
  \item $(x+3)(x-3)=x^2-9$.
  \item $(2x-1)^2=4x^2-4x+1$.
  \item $x^2-9=(x-3)(x+3)$.
  \item $4x^2-4x+1=(2x-1)^2$.
  \item $(x+2)(x^2-2x+4)=x^3+8$.
  \item $x^3+8=(x+2)(x^2-2x+4)$.
\end{enumerate}

%────────────────────────────────────────────────────────────────
\solutiontitle{4}{27}{R\'esoudre avec le discriminant (suite)}

\begin{enumerate}
  \item $x(x-\sqrt{5})=0$~: $x=0$ ou $x=\sqrt{5}$.
  \item $x^2+x-6=(x-2)(x+3)=0$~: $x=2$ ou $x=-3$.
  \item $(x-1)^2=0$~: racine double $x=1$.
  \item $x^2-2x+1=(x-1)^2=0$~: racine double $x=1$.
\end{enumerate}

%────────────────────────────────────────────────────────────────
\solutiontitle{4}{28}{Probl\`eme de vitesse}

Soit $v$ la vitesse actuelle (km/h).
$\frac{360}{v}-\frac{360}{v+30}=1\Rightarrow v^2+30v-10800=0$~; $\Delta=900+43200=44100=210^2$.
$v=\frac{-30+210}{2}=90$~km/h.\hfill$\square$

%────────────────────────────────────────────────────────────────
\solutiontitle{4}{29}{Signe d'un produit de facteurs}

\begin{enumerate}
  \item $x^2(x-1)$~: $x^2\geq0$, signe de $(x-1)$.
        $<0$ sur $]-\infty,0[\cup]0,1[$, $=0$ en $0$ et $1$, $>0$ sur $]1,+\infty[$.
  \item $(x+2)^2(x-3)$~: nul en $-2$. Signe de $(x-3)$~: n\'egatif avant $3$.
        $\leq0$ sur $]-\infty,3]$~; $=0$ en $-2$ et $3$~; $>0$ sur $]3,+\infty[$.
  \item $\frac{(x-3)(x+3)}{(x+1)^2}$~: $(x+1)^2>0$ pour $x\ne-1$.
        Signe = signe de $(x-3)(x+3)$~: $+$ sur $]-\infty,-3[$, $-$ sur $]-3,3[$ ($x\ne-1$),
        $+$ sur $]3,+\infty[$.
  \item $(x-2)(x+2)(x-1)(x+1)$~: z\'eros $\pm1,\pm2$.
        Signe~: $+,-,+,-,+$ sur les cinq intervalles.
\end{enumerate}

%────────────────────────────────────────────────────────────────
\solutiontitle{4}{30}{Valeurs interdites et simplification}

\begin{enumerate}
  \item $\frac{x(x-1)}{x-1}=x$ ($x\ne1$). $D=\R\setminus\{1\}$.
  \item $\frac{x(x^2-4)}{x^2-4}=x$ ($x\ne\pm2$). $D=\R\setminus\{-2,2\}$.
  \item $\frac{(x+1)(x+2)}{(x+1)(x-1)}=\frac{x+2}{x-1}$ ($x\ne\pm1$). $D=\R\setminus\{-1,1\}$.
  \item $\frac{x^2(x+1)}{x(x+1)}=x$ ($x\ne0,-1$). $D=\R\setminus\{-1,0\}$.
\end{enumerate}

%────────────────────────────────────────────────────────────────
\solutiontitle{4}{31}{\logicoalgo\ Algorithme~: r\'esoudre $ax^2+bx+c=0$}

\begin{enumerate}
  \item Les trois cas $\Delta<0$, $\Delta=0$, $\Delta>0$ sont disjoints et exhaustifs
        (tout r\'eel v\'erifie exactement l'un d'eux). La structure couvre tous les cas.\checkmark
  \item Traces~:
  \begin{itemize}
    \item $a=1,b=-5,c=6$~: $\Delta=1>0$~; deux racines $x_1=2$, $x_2=3$.
    \item $a=1,b=2,c=1$~: $\Delta=0$~; racine double $x_0=-1$.
    \item $a=1,b=0,c=1$~: $\Delta=-4<0$~; pas de racine r\'eelle.
  \end{itemize}
  \item Ajouter avant le calcul de $\Delta$~: \textbf{Si} $a=0$, traiter le cas lin\'eaire
        $bx+c=0$ (si $b\ne0$~: $x=-c/b$~; si $b=c=0$~: tout r\'eel~; sinon impossible).
  \item
\begin{lstlisting}[language=Python]
import math
def resoudre(a, b, c):
    if a == 0:
        if b != 0: print(f"Lineaire: x = {-c/b}")
        elif c == 0: print("Tout reel")
        else: print("Impossible")
        return
    d = b**2 - 4*a*c
    if d < 0: print("Pas de racine reelle")
    elif d == 0: print(f"Racine double: {-b/(2*a)}")
    else:
        print(f"x1={(-b-math.sqrt(d))/(2*a):.4f}")
        print(f"x2={(-b+math.sqrt(d))/(2*a):.4f}")
\end{lstlisting}
\end{enumerate}

%────────────────────────────────────────────────────────────────
\solutiontitle{4}{32}{Signe strict d'un trin\^ome}

\begin{enumerate}
  \item $x^2+x+1=\bigl(x+\tfrac{1}{2}\bigr)^2+\tfrac{3}{4}\geq\tfrac{3}{4}>0$
        pour tout $x\in\R$.\hfill$\square$
  \item $x^4+x^2+1\geq1>0$ pour tout $x\in\R$.\hfill$\square$
  \item $x^2-x+1=\bigl(x-\tfrac{1}{2}\bigr)^2+\tfrac{3}{4}>0$, la division est licite.
        $3(x^2-x+1)-(x^2+x+1)=2(x-1)^2\geq0$~:
        donc $\frac{x^2+x+1}{x^2-x+1}\leq3$ (\'egalit\'e en $x=1$).\hfill$\square$
\end{enumerate}

%────────────────────────────────────────────────────────────────
\solutiontitle{4}{33}{Racines et coefficients~: Vi\`ete appliqu\'e}

$x_1+x_2=p$, $x_1x_2=q$.
\begin{enumerate}
  \item $x_1^2+x_2^2=(x_1+x_2)^2-2x_1x_2=p^2-2q$.
  \item $(x_1-x_2)^2=p^2-4q=\Delta$.
  \item $x_1^3+x_2^3=(x_1+x_2)(x_1^2-x_1x_2+x_2^2)=p(p^2-3q)$.
  \item $p=3$, $q=1$~: $x_1^3+x_2^3=3(9-3)=18$.\checkmark
\end{enumerate}

%────────────────────────────────────────────────────────────────
\solutiontitle{4}{34}{Fraction rationnelle~: comportement}

$F(x)=\dfrac{x^2-1}{x^2+1}$.
\begin{enumerate}
  \item $D=\R$ (le d\'enominateur $x^2+1>0$ toujours).
  \item $F(x)<1$~: $x^2-1<x^2+1$, vrai~; $F(x)\geq-1$~: $2x^2\geq0$, vrai.
        Donc $-1\leq F(x)<1$ ($F(0)=-1$, borne atteinte).
  \item $F(-x)=\dfrac{x^2-1}{x^2+1}=F(x)$~: $F$ est \textbf{paire}.\hfill$\square$
  \item $F(x)\geq0\Leftrightarrow x^2\geq1$~: $F>0$ sur $]-\infty,-1[\cup]1,+\infty[$,
        $F=0$ en $\pm1$, $F<0$ sur $]-1,1[$.
\end{enumerate}

%────────────────────────────────────────────────────────────────
\solutiontitle{4}{35}{Polyn\^ome de degr\'e 3 \`a racine double}

\begin{enumerate}
  \item $P'(x)=3x^2-3=3(x-1)(x+1)$. $P(1)=1-3+2=0$~: racine double $r=1$.
        $P(-1)=-1+3+2=4\ne0$~: $-1$ n'est pas racine de $P$.
  \item $P(x)=(x-1)^2(x-r_3)$~; terme constant~: $(1)(-r_3)=2\Rightarrow r_3=-2$.
        $P(x)=(x-1)^2(x+2)$. V\'erif.~: $(x^2-2x+1)(x+2)=x^3-3x+2$.\checkmark
  \item $Q'(x)=3x^2-3$~; $Q(-1)=-1+3-2=0$~: racine double $r=-1$.
        $Q(x)=(x+1)^2(x-2)$. V\'erif.~: $(x^2+2x+1)(x-2)=x^3-3x-2$.\checkmark\hfill$\square$
\end{enumerate}

%────────────────────────────────────────────────────────────────
\solutiontitle{4}{36}{Probl\`eme \`a param\`etre et trin\^ome}

D\'eterminant~: $\det=k^2-1=(k-1)(k+1)$.
Unique solution $\Leftrightarrow k\ne\pm1$.

Pour $k\ne\pm1$~: par combinaison lin\'eaire, $x=y=\dfrac{1}{k+1}$.

Si $k=1$~: \'equations identiques~; infinit\'e de solutions.
Si $k=-1$~: $-x+y=1$ et $x-y=1$~: incompatible, aucune solution.

%────────────────────────────────────────────────────────────────
\solutiontitle{4}{37}{Localisation de racines}

\begin{enumerate}
  \item Th\'eor\`eme des valeurs interm\'ediaires (admis)~: $P$ continue,
        $P(a)<0<P(b)$, donc $\exists c\in]a,b[,\;P(c)=0$.\hfill$\square$
  \item $P(0)=1>0$, $P(1)=-1<0$~: racine dans $]0,1[$.\checkmark\quad
        $P(1)=-1<0$, $P(2)=3>0$~: racine dans $]1,2[$.\checkmark
  \item Affinage~: $P(0{,}3)>0$, $P(0{,}4)<0$~: racine dans $]0{,}3;0{,}4[$.
        $P(1{,}5)<0$, $P(1{,}6)>0$~: racine dans $]1{,}5;1{,}6[$.\hfill$\square$
\end{enumerate}

%────────────────────────────────────────────────────────────────
\solutiontitle{4}{38}{\logiq\ Nier une \'equation ou in\'equation}

\begin{enumerate}
  \item $\forall x\in\R,\;x^2-3x+2\ne0$.
  \item $\exists x\in\R,\;x^2+1\leq0$.
  \item $\exists x\in[0,4],\;2x^2-5x+3<0$.
  \item \og $ax^2+bx+c=0$ n'admet aucune solution r\'eelle.\fg
\end{enumerate}

%────────────────────────────────────────────────────────────────
\solutiontitle{4}{39}{\logiq\ Conditions sur le discriminant}

\begin{enumerate}
  \item $\exists x\in\R,\;f(x)=0\Leftrightarrow\Delta\geq0$ (pour $a\ne0$).
  \item \textbf{Fausse.} Si $\Delta=0$, il y a exactement une racine (double).
        \og Deux racines\fg\ requiert $\Delta>0$.
  \item N\'egation~: $\exists x\in\R,\;f(x)\leq0$.
        Condition~: $a\leq0$ ou $\Delta\geq0$
        (positivit\'e universelle exige $a>0$ et $\Delta<0$).
\end{enumerate}

%────────────────────────────────────────────────────────────────
\solutiontitle{4}{40}{\logiq\ Factorisation et logique}

\begin{enumerate}
  \item Par le th\'eor\`eme de B\'ezout~:
        $P(r)=0\Leftrightarrow\exists Q(x),\;P(x)=(x-r)Q(x)$.
  \item \textbf{Vraie.} $x^2+1>0$ toujours ($\Delta=-4<0$)~:
        la seule racine r\'eelle est $x=1$.
  \item $x^2\geq0$ donc $x^2+1\geq1>0$~: raisonnement \textbf{universel direct}.\hfill$\square$
\end{enumerate}

%────────────────────────────────────────────────────────────────
\solutiontitle{4}{41}{\logiq\ Discussion selon un param\`etre}

$\Delta=m^2-4$.
\begin{enumerate}
  \item $|m|>2$~: deux solutions~; $m=\pm2$~: une solution double~; $|m|<2$~: aucune.
  \item \og L'\'equation a deux solutions si et seulement si $|m|>2$.\fg
  \item N\'egation de deux solutions~: $|m|\leq2$.
\end{enumerate}

%────────────────────────────────────────────────────────────────
\solutiontitle{4}{42}{\logiq\ Raisonnement dans les syst\`emes}

\begin{enumerate}
  \item Le syst\`eme $ax+by=e$, $cx+dy=f$ a une unique solution
        $\Leftrightarrow ad-bc\ne0$.
  \item Suppos. par l'absurde~: $D=0$ et unique solution $(x_0,y_0)$.
        La deuxi\`eme ligne est proportionnelle \`a la premi\`ere~;
        il y a soit incompatibilit\'e soit infinit\'e de solutions.
        Contradiction.\hfill$\square$
  \item $D=0$~: lignes parall\`eles distinctes ($\Rightarrow0$ solution)
        ou confondues ($\Rightarrow\infty$ solutions).
\end{enumerate}

%────────────────────────────────────────────────────────────────
\solutiontitle{4}{43}{\logiq\ Fractions rationnelles et domaine}

\begin{enumerate}
  \item $D_f=\{x\in\R\mid x\ne3\}$.
  \item $D_f=\{x\in\R\mid x\ne\pm2\}$.
  \item $D_f=\R$ (car $\forall x,\;x^2+1>0$).
  \item $D_f=\{x\in\R\mid x\geq0\wedge x\ne1\}$.
\end{enumerate}

%────────────────────────────────────────────────────────────────
\solutiontitle{4}{44}{\algo\ Algorithme~: r\'esoudre $ax^2+bx+c=0$}

\begin{enumerate}
  \item Ajouter en t\^ete~: \textbf{Si} $a=0$, traiter $bx+c=0$ puis sortir.
  \item Apr\`es calcul de $x_1$~: calculer $ax_1^2+bx_1+c$ et afficher si nul.
  \item
\begin{lstlisting}[language=Python]
import math
def solve(a, b, c):
    if a == 0:
        return [-c/b] if b != 0 else []
    d = b*b - 4*a*c
    if d < 0: return []
    if d == 0: return [-b/(2*a)]
    return [(-b-math.sqrt(d))/(2*a), (-b+math.sqrt(d))/(2*a)]
\end{lstlisting}
\end{enumerate}

%────────────────────────────────────────────────────────────────
\solutiontitle{4}{45}{\algo\ Algorithme~: signe d'un trin\^ome en un point}

\begin{enumerate}
  \item $f(x)=x^2-5x+6$~: $f(1)=2$ (Positif), $f(2)=0$ (Nul),
        $f(3)=0$ (Nul), $f(4)=2$ (Positif).
  \item Ajouter la ligne \og Afficher $v$\fg\ avant le test.
  \item
\begin{lstlisting}[language=Python]
a, b, c = 1, -5, 6
for x in range(-10, 11):
    if a*x**2 + b*x + c < 0:
        print(x)
\end{lstlisting}
\end{enumerate}

%────────────────────────────────────────────────────────────────
\solutiontitle{4}{46}{\algo\ Algorithme~: tableau de signe d'un polyn\^ome de degr\'e 3}

$P(x)=(x-1)(x-2)(x-3)$.

\begin{center}
\begin{tabular}{c|ccccccc}
$x$ & $-1$ & $0$ & $1$ & $2$ & $3$ & $4$ & $5$ \\\hline
$P(x)$ & $-24$ & $-6$ & $0$ & $0$ & $0$ & $6$ & $24$ \\
Signe & $-$ & $-$ & $0$ & $0$ & $0$ & $+$ & $+$
\end{tabular}
\end{center}

\begin{lstlisting}[language=Python]
P = lambda x: (x-1)*(x-2)*(x-3)
for x in range(-1, 6):
    print(x, P(x))
\end{lstlisting}

%────────────────────────────────────────────────────────────────
\solutiontitle{4}{47}{\algo\ Algorithme~: trouver les racines par balayage \computer}

On parcourt $[-100,100]$ par pas $0{,}01$ et on d\'etecte les changements de signe.

\begin{lstlisting}[language=Python]
a, b, c = 1, -5, 6
x, prev = -100.0, None
while x <= 100.0:
    curr = a*x**2 + b*x + c
    if prev is not None and prev * curr < 0:
        print(f"Racine ~ {x:.2f}")
    prev, x = curr, round(x + 0.01, 10)
\end{lstlisting}
Ce balayage localise deux racines proches de $x=2{,}00$ et $x=3{,}00$.

%────────────────────────────────────────────────────────────────
\solutiontitle{4}{48}{\algo\ Algorithme~: sch\'ema de H\"orner appliqu\'e}

\begin{lstlisting}[language=Python]
def horner(coeffs, x0):
    r = 0
    for c in coeffs:
        r = r * x0 + c
    return r

coeffs = [1, -6, 11, -6]
for x0 in [1, 2, 3, 4]:
    print(f"P({x0}) = {horner(coeffs, x0)}")
# Resultats: 0, 0, 0, 6
\end{lstlisting}

%────────────────────────────────────────────────────────────────
\solutiontitle{4}{49}{\algo\ Algorithme~: r\'esoudre une in\'equation par balayage}

$P(x)=x^3-3x=x(x-\sqrt{3})(x+\sqrt{3})$. R\'esolution alg\'ebrique~:
$P>0$ sur $]-\sqrt{3},0[\cup]\sqrt{3},+\infty[$.
Entiers solutions dans $\{-3,\ldots,3\}$~: $x\in\{-1,2,3\}$.

\begin{lstlisting}[language=Python]
for x in range(-3, 4):
    if x**3 - 3*x > 0:
        print(x)  # Affiche: -1, 2, 3
\end{lstlisting}

%────────────────────────────────────────────────────────────────
\subsection*{Probl\`emes}
\addcontentsline{toc}{subsection}{Probl\`emes}

%────────────────────────────────────────────────────────────────
\psolutiontitle{4}{1}{Applications g\'eom\'etriques des \'equations}

\textbf{(1)} Entiers cons\'ecutifs $n$ et $n+1$~: $(n+1)^2-n^2=2n+1=47\Rightarrow n=23$.
Les entiers sont $\mathbf{23}$ et $\mathbf{24}$.

\textbf{(2)} Terrain $a\times b$, all\'ee de $2$~m~: dimensions totales $(a+4)\times(b+4)$.
$ab=24$ et $(a+4)(b+4)=80\Rightarrow4(a+b)=40\Rightarrow a+b=10$.
L'\'equation $t^2-10t+24=0$ a pour racines $4$ et $6$.
Les dimensions du terrain sont donc $\mathbf{4\ \text{m}\times6\ \text{m}}$.

\textbf{(3)} $\frac{1}{t}+\frac{1}{t+5}=\frac{1}{6}\Rightarrow t^2-7t-30=0$~;
$\Delta=169=13^2$~; $t=10$~min. Robinet~A~: $10$~min~; B~: $15$~min.

%────────────────────────────────────────────────────────────────
\psolutiontitle{4}{2}{Intersection droite--parabole}

$\mathcal{P}$~: $y=x^2-2x-3$~; $d_k$~: $y=kx+1$.

\textbf{(1)} $x^2-2x-3=1\Rightarrow x^2-2x-4=0$~; $\Delta=20$~;
intersections $(1+\sqrt{5},1)$ et $(1-\sqrt{5},1)$.

\textbf{(2)} $x^2-(2+k)x-4=0$~; $\Delta=(2+k)^2+16>0$ toujours.
La droite est \textbf{toujours s\'ecante}~; aucune tangente de la forme $y=kx+1$.

\textbf{(3)} S\'ecante en deux points pour tout $k\in\R$ (voir (2)).

\textbf{(4)} $k=3$~: $x^2-5x-4=0$~; $\Delta=41$.
Distance~: $|x_1-x_2|\sqrt{1+k^2}=\sqrt{41}\cdot\sqrt{10}=\sqrt{410}$.\hfill$\square$

%────────────────────────────────────────────────────────────────
\psolutiontitle{4}{3}{Fractions rationnelles et physique}

$v=\dfrac{3(x-2)(x+2)}{(x-2)^2}=\dfrac{3(x+2)}{x-2}$ pour $x\ne2$.

\textbf{(1)} $D=\R\setminus\{2\}$.
\textbf{(2)} $v=\dfrac{3(x+2)}{x-2}$.
\textbf{(3)} $v>3\Rightarrow\dfrac{4}{x-2}>0\Rightarrow x>2$.
\textbf{(4)} $v=0\Rightarrow x=-2$~: vitesse nulle pour $x=-2$.

%────────────────────────────────────────────────────────────────
\psolutiontitle{4}{4}{Trin\^ome du second degr\'e et extremum}

\textbf{(1)} $f(x)=a(x-\alpha)^2+\beta$ avec $\alpha=-\dfrac{b}{2a}$, $\beta=-\dfrac{\Delta}{4a}$.\hfill$\square$

\textbf{(2)} $a>0$~: $a(x-\alpha)^2\geq0\Rightarrow f\geq\beta$ (minimum).
$a<0$~: $f\leq\beta$ (maximum).\hfill$\square$

\textbf{(3)} $h(t)=-5t^2+20t+1$~: $\alpha=2$~s, $\beta=21$~m.
Hauteur maximale~: $\mathbf{21\,\text{m}}$ atteinte \`a $t=2$~s.

\textbf{(4)} $h(t)=0$~: $\Delta=400+20=420$~;
$t=2+\tfrac{\sqrt{105}}{5}\approx4{,}05$~s.

%────────────────────────────────────────────────────────────────
\psolutiontitle{4}{5}{Polyn\^ome de degr\'e 3 et factorisation}

\textbf{(1)} $a(a-1)(a+1)$~: produit de trois entiers cons\'ecutifs.
L'un est pair ($\Rightarrow2\mid P$) et l'un est divisible par $3$ ($\Rightarrow3\mid P$).
Donc $6\mid P(a)$.\hfill$\square$

\textbf{(2)} $x^3-y^3$~: $P(y)=0$~; division euclidienne~:
$x^3-y^3=(x-y)(x^2+xy+y^2)$.\hfill$\square$

\textbf{(3)} $x^3-1=(x-1)(x^2+x+1)$.\hfill$\square$

\textbf{(4)} $\Delta(x^2+x+1)=-3<0$~: $x^2+x+1>0$ toujours.
Signe de $x^3-1$~: celui de $(x-1)$.
$<0$ pour $x<1$~; $=0$ en $x=1$~; $>0$ pour $x>1$.

%─────────────────────────────────────────────────────────��──────
\psolutiontitle{4}{6}{Syst\`eme et g\'eom\'etrie}

\textbf{(1)} $2x-1=-x+5\Rightarrow x=2$, $y=3$.
Point~$A=(2,3)$.

\textbf{(2)} $D_3$~: $y=3x+1$.

\textbf{(3)} $D_1\cap D_3$~: $2x-1=3x+1$, donc $B=(-2,-5)$.
$D_2\cap D_3$~: $-x+5=3x+1$, donc $C=(1,4)$.
Les sommets sont $A=(2,3)$, $B=(-2,-5)$ et $C=(1,4)$.

\textbf{(4)} $\vv{AB}=(-4,-8)$ et $\vv{AC}=(-1,1)$.
L'aire vaut $\frac12\left|(-4)\times1-(-8)\times(-1)\right|=\mathbf{6}$.

%────────────────────────────────────────────────────────────────
\psolutiontitle{4}{7}{D\'enombrement de racines positives}

\textbf{(1)} Tester $x=2$~: $4-2(m+2)+2m=0$. Donc $2$ est toujours racine.
Autre racine~: $x_2=m$ (par Vi\`ete).
Deux racines $>0$~: $m>0$.

\textbf{(2)} $P(x)=x(x-1)(x-2)$~: racines $0$, $1$, $2$.
Racines strictement positives~: $1$ et $2$.

\textbf{(3)} $Q(x)=(x+1)(x^2+1)$~: unique racine r\'eelle $x=-1$.\hfill$\square$

%────────────────────────────────────────────────────────────────
\psolutiontitle{4}{8}{\'Etude compl\`ete d'une fraction rationnelle}

$F(x)=\dfrac{(x-1)(x-3)}{(x-1)(x+1)}=\dfrac{x-3}{x+1}$ pour $x\ne1,-1$.

\textbf{(1)} $D=\R\setminus\{-1,1\}$.
\textbf{(2--3)} Factorisation et simplification ci-dessus.
\textbf{(4)} Tableau~: z\'ero en $3$, exclu en $-1$~:
$F>0$ sur $]-\infty,-1[\cup]3,+\infty[$~; $F<0$ sur $]-1,1[\cup]1,3[$.
\textbf{(5)} $F\leq0$~: $x\in]-1,1[\cup]1,3]$.

%────────────────────────────────────────────────────────────────
\psolutiontitle{4}{9}{Sch\'ema de H\^orner pour les racines}

\textbf{(1)} $P(x)=x^3-3x^2+x-5$, $r=2$. Coefficients~: $[1,-3,1,-5]$.
$1\to-1\to-1\to-7$. $P(2)=-7$.

\textbf{(2)} $P(x)=x^3-6x^2+11x-6$, $r=1$. Coefficients~: $[1,-6,11,-6]$.
$1\to-5\to6\to0$. Reste $=P(1)=0$.\checkmark\quad Quotient~: $x^2-5x+6$.

\textbf{(3)} $(x-1)(x^2-5x+6)=x^3-6x^2+11x-6$.\checkmark

%────────────────────────────────────────────────────────────────
\psolutiontitle{4}{10}{Probl\`eme d'optimisation}

\textbf{(1)} $V(x)=x(20-2x)^2=4x(10-x)^2$.
\textbf{(2)} $x\in]0,10[$.
\textbf{(3)} $V(1)=324$, $V(2)=512$, $V(3)=588$, $V(4)=576$~cm$^3$.
\textbf{(4)} $V'(x)=4(3x-10)(x-10)$~; $V'=0$ en $x=\tfrac{10}{3}$ et $x=10$.
$V'>0$ sur $]0,\tfrac{10}{3}[$~; maximum en $x=\tfrac{10}{3}$.
$V\!\bigl(\tfrac{10}{3}\bigr)=\tfrac{16000}{27}\approx592{,}6$~cm$^3$.\hfill$\square$

%────────────────────────────────────────────────────────────────
\psolutiontitle{4}{11}{Propri\'et\'e des racines -- suite}

$P(x)=(x-r_1)(x-r_2)(x-r_3)$ identifi\'e avec $x^3+px+q$~:
\begin{enumerate}
  \item $r_1+r_2+r_3=0$.\hfill$\square$
  \item $r_1r_2+r_1r_3+r_2r_3=p$.\hfill$\square$
  \item $r_1r_2r_3=-q$.\hfill$\square$
  \item $r_1=1$~: $r_2+r_3=-1$, $r_2r_3=-2$~;
        $t^2+t-2=(t+2)(t-1)=0$~; $r_2=-2$, $r_3=1$ (racine double).
        $P(x)=(x-1)^2(x+2)$.\checkmark
\end{enumerate}

%────────────────────────────────────────────────────────────────
\psolutiontitle{4}{12}{Double in\'egalit\'e}

\textbf{(1)} $0<\dfrac{x+1}{x-2}\leq3$.
$\frac{x+1}{x-2}>0$~: $x<-1$ ou $x>2$.
$\frac{x+1}{x-2}\leq3\Leftrightarrow\frac{-2x+7}{x-2}\geq0$~:
pour $x>2$~: $x\leq\frac{7}{2}$~; pour $x<-1$~: toujours vrai.
\textbf{Solution~:} $]-\infty,-1[\cup]2,\tfrac{7}{2}]$.

\textbf{(2)} Partie droite~: $x^2-4\leq x^2+1$~: toujours.
Partie gauche~: $2x^2\geq3\Rightarrow|x|\geq\tfrac{\sqrt{6}}{2}$.
\textbf{Solution~:} $\bigl]-\infty,-\tfrac{\sqrt{6}}{2}\bigr]\cup\bigl[\tfrac{\sqrt{6}}{2},+\infty\bigr[$.

\textbf{(3)} Cas $x>0$~: $\frac{1}{x}<x\Rightarrow x>1$~;
$x<\frac{1}{x^2}\Rightarrow x^3<1\Rightarrow x<1$~: contradiction.
Cas $x<0$~: $\frac{1}{x}<x$ (mult. par $x^2>0$)~: $x<x^3\Rightarrow x^3-x>0\Rightarrow x(x^2-1)>0$.
Pour $-1<x<0$~: $x<0$ et $x^2-1<0$~: produit $>0$.\checkmark
$x<\frac{1}{x^2}$~: $x^3<1$~: vrai pour $x<0$.
\textbf{Solution~:} $]-1,0[$.

%────────────────────────────────────────────────────────────────
\psolutiontitle{4}{13}{\logiq\ Logique et \'equations param\'etriques}

\textbf{(1)} Trin\^ome ssi $m-1\ne0$, i.e. $m\ne1$.
\textbf{(2)} $\Delta=12-8m$~; deux racines distinctes ssi $\Delta>0\Leftrightarrow m<\tfrac{3}{2}$.
$P\Leftrightarrow Q$ par d\'efinition du discriminant.\hfill$\square$
\textbf{(3)} Vi\`ete~: $x_1+x_2=\frac{2m}{m-1}>0$, $x_1x_2=\frac{m+3}{m-1}>0$, $\Delta\geq0$.
Intersection~: $m<-3$ ou $1<m\leq\tfrac{3}{2}$.

%────────────────────────────────────────────────────────────────
\psolutiontitle{4}{14}{\logiq\ Syst\`emes et logique}

\textbf{(1)} La deuxi\`eme \'equation vaut $2\times$ la premi\`ere.
Si $b=2a$~: infinit\'e de solutions~; si $b\ne2a$~: syst\`eme incompatible.
\textbf{(2)} G\'en\'eral~: $D=0\Rightarrow$ soit $0$ solution (lignes parall\`eles distinctes),
soit $\infty$ solutions (lignes confondues).

%────────────────────────────────────────────────────────────────
\psolutiontitle{4}{15}{\algo\ Programme Python~: r\'esolution compl\`ete}

\begin{lstlisting}[language=Python]
import math

def resoudre(a, b, c):
    if a == 0:
        if b == 0:
            print("Tout reel" if c == 0 else "Impossible")
        else:
            x = -c / b
            print(f"Lineaire: x = {x}")
            print(f"Verif: {b*x + c}")
        return
    delta = b**2 - 4*a*c
    if delta < 0:
        print("Pas de racine reelle")
    elif delta == 0:
        x0 = -b / (2*a)
        print(f"Racine double: {x0}")
        print(f"Verif: {round(a*x0**2+b*x0+c, 10)}")
    else:
        x1 = (-b - math.sqrt(delta)) / (2*a)
        x2 = (-b + math.sqrt(delta)) / (2*a)
        print(f"x1={x1:.6f}, x2={x2:.6f}")
        print(f"Verif: {round(a*x1**2+b*x1+c,8)}, {round(a*x2**2+b*x2+c,8)}")
\end{lstlisting}

%────────────────────────────────────────────────────────────────
\psolutiontitle{4}{16}{\algo\ Division euclidienne par H\"orner}

\begin{lstlisting}[language=Python]
def division_horner(coeffs, r):
    q = [coeffs[0]]
    for i in range(1, len(coeffs) - 1):
        q.append(q[-1] * r + coeffs[i])
    reste = q[-1] * r + coeffs[-1]
    return q, reste

q, reste = division_horner([1, -6, 11, -6], 1)
print("Quotient:", q)   # [1, -5, 6]
print("Reste:", reste)  # 0  (= P(1))
\end{lstlisting}
Th\'eor\`eme de B\'ezout~: le reste $= P(1) = 0$.\checkmark

%────────────────────────────────────────────────────────────────
\subsection*{D\'efis}
\addcontentsline{toc}{subsection}{D\'efis}

%────────────────────────────────────────────────────────────────
\noindent\phantomsection
{\color{BO}\bfseries\small D\'efi~1~\textendash\ \textit{\'Equation bicarr\'ee}}
\par\smallskip

Poser $u=x^2$~: $u^2-5u+4=(u-1)(u-4)=0\Rightarrow u\in\{1,4\}$.
$u=1\Rightarrow x=\pm1$~; $u=4\Rightarrow x=\pm2$.
\textbf{Solutions~: $\{-2,-1,1,2\}$.}

$x^4-5x^2+4\leq0\Leftrightarrow(u-1)(u-4)\leq0\Leftrightarrow1\leq u\leq4\Leftrightarrow1\leq x^2\leq4$.
\textbf{Solution~: $[-2,-1]\cup[1,2]$.}

%────────────────────────────────────────────────────────────────
\noindent\phantomsection
{\color{BO}\bfseries\small D\'efi~2~\textendash\ \textit{Irrationalit\'e et trin\^ome}}
\par\smallskip

\textbf{(1)} Racines de $x^2-3$~: $\pm\sqrt{3}$ ($\Delta=12$).

\textbf{(2)} Supposons $\sqrt{3}=p/q$ irr\'eductible.
$p^2=3q^2\Rightarrow3\mid p\Rightarrow p=3k\Rightarrow9k^2=3q^2\Rightarrow3\mid q$.
Contradiction avec $\pgcd(p,q)=1$.\hfill$\square$

\textbf{(3)} Pour $P(x)=x^2-n$, $n$ non carr\'e parfait~:
si $\sqrt{n}=p/q$, alors $p^2=nq^2$. Par un argument identique de divisibilit\'e,
tout facteur premier de $n$ divise \`a la fois $p$ et $q$, contredisant l'irr\'eductibilit\'e.\hfill$\square$

%────────────────────────────────────────────────────────────────
\noindent\phantomsection
{\color{BO}\bfseries\small D\'efi~3~\textendash\ \textit{Triplets de Pythagore}}
\par\smallskip

\textbf{(1)} $a^2+b^2=(m^2-n^2)^2+(2mn)^2=m^4-2m^2n^2+n^4+4m^2n^2=(m^2+n^2)^2=c^2$.\hfill$\square$

\textbf{(2)} Triplets~:
$(2,1)\to(3,4,5)$~; $(3,2)\to(5,12,13)$~;
$(4,1)\to(15,8,17)$~; $(4,3)\to(7,24,25)$.

\textbf{(3)} $c=m^2+n^2$ est toujours un entier. Mais tout entier n'est pas hypot\'enuse~:
par exemple $3$ n'est hypot\'enuse d'aucun triplet primitif.

\textbf{(4)} $5^2+12^2=25+144=169=13^2$.\checkmark\quad
$8^2+15^2=64+225=289=17^2$.\checkmark

%────────────────────────────────────────────────────────────────
\noindent\phantomsection
{\color{BO}\bfseries\small D\'efi~4~\textendash\ \textit{Somme et produit de trois racines}}
\par\smallskip

\textbf{(1)} $(x-r_1)(x-r_2)(x-r_3)=x^3-(r_1+r_2+r_3)x^2
+(r_1r_2+r_1r_3+r_2r_3)x-r_1r_2r_3$.
Identification avec $x^3+bx^2+cx+d$~:
$r_1+r_2+r_3=-b$, $r_1r_2+r_1r_3+r_2r_3=c$, $r_1r_2r_3=-d$.\hfill$\square$

\textbf{(2)} Racines $1,2,3$~: $b=-(1+2+3)=-6$, $c=2+3+6=11$, $d=-6$.
Polyn\^ome~: $x^3-6x^2+11x-6$.

\textbf{(3)} Racines $\sqrt{2},-\sqrt{2},1$~: $b=-(1+0)=-1$, $c=(-2+\sqrt{2}-\sqrt{2})=-2$,
$d=-(-2)=2$. Polyn\^ome~: $x^3-x^2-2x+2=(x-1)(x^2-2)$.

%────────────────────────────────────────────────────────────────
\noindent\phantomsection
{\color{BO}\bfseries\small D\'efi~5~\textendash\ \textit{M\'ethode de Cardano~: introduction}}
\par\smallskip

\textbf{(1)} $(u+v)^3=u^3+3u^2v+3uv^2+v^3=u^3+v^3+3uv(u+v)$.
Avec $x=u+v$~: $x^3=u^3+v^3+3uvx$.
Pour $x^3+px+q=0$~: il suffit d'avoir $u^3+v^3=-q$ et $3uv=-p$, soit $uv=-p/3$.

\textbf{(2)} $u^3$ et $v^3$ v\'erifient~: somme $=-q$, produit $=-p^3/27$.
\'Equation du second degr\'e~: $t^2+qt-p^3/27=0$.

\textbf{(3)} $x^3-3x+2=0$~: $p=-3$, $q=2$.
\'Equation~: $t^2+2t+1=(t+1)^2=0\Rightarrow u^3=v^3=-1$.
$u=v=-1\Rightarrow x=u+v=-2$, ou prendre les racines cubiques complexes pour retrouver $x=1$.\hfill$\square$

%────────────────────────────────────────────────────────────────
\noindent\phantomsection
{\color{BO}\bfseries\small D\'efi~6~\textendash\ \textit{R\`egle des signes de Descartes}}
\par\smallskip

\textbf{(1)} $P(x)=x^3-x^2-x+1$~: coefficients $+,-,-,+$.
Changements de signe~: $+\to-$ (oui), $-\to-$ (non), $-\to+$ (oui).
Deux changements~; au plus \textbf{deux} racines positives.

\textbf{(2)} $P(x)=x^2(x-1)-(x-1)=(x-1)(x^2-1)=(x-1)^2(x+1)$.
Racines~: $x=1$ (double) et $x=-1$. Racines positives~: $x=1$ (une seule). Coh\'erent (2 est un majorant).\checkmark

\textbf{(3)} $Q(x)=x^4-x^3+x-1$~: coefficients $+,-,0,+,-$.
Changements~: $+\to-$, $-\to+$, $+\to-$~: trois changements, donc $\leq3$ racines positives.
$Q(x)=x^3(x-1)+(x-1)=(x-1)(x^3+1)=(x-1)(x+1)(x^2-x+1)$.
Racines r\'eelles~: $x=1$ et $x=-1$. Racine positive~: $x=1$.\checkmark

\textbf{(4)} $P(-x)=-(x^3+x^2-x-1)=-(x+1)^2(x-1)$.
Pour $P(-y)$ avec $y=-x$, un changement de signe~: $\leq1$ racine n\'egative.
V\'erif.~: racine n\'egative de $P$~: $x=-1$.\checkmark

%────────────────────────────────────────────────────────────────
\noindent\phantomsection
{\color{BO}\bfseries\small D\'efi~7~\textendash\ \textit{\logiq\ Vi\`ete et logique des racines}}
\par\smallskip

\textbf{(1)} Si $r_1,r_2$ sont les racines de $x^2+px+q=0$~:
$r_1+r_2=-p$ et $r_1r_2=q$.

\textbf{(2)} R\'eciproque~: si $r_1+r_2=-p$ et $r_1r_2=q$, alors
$r_1$ et $r_2$ sont racines de $t^2-(r_1+r_2)t+r_1r_2=t^2+pt+q=0$.\hfill$\square$

\textbf{(3)} $r_1,r_2>0\Leftrightarrow r_1+r_2>0\wedge r_1r_2>0\wedge\Delta\geq0$
$\Leftrightarrow p<0\wedge q>0\wedge\Delta\geq0$.\hfill$\square$

%────────────────────────────────────────────────────────────────
\noindent\phantomsection
{\color{BO}\bfseries\small D\'efi~8~\textendash\ \textit{\algo\ M\'ethode de Newton \computer}}
\par\smallskip

\textbf{(1)} $f(x)=x^2-a$, $f'(x)=2x$.
$u_{n+1}=u_n-\frac{u_n^2-a}{2u_n}=\frac{u_n^2+a}{2u_n}=\frac{1}{2}\!\left(u_n+\frac{a}{u_n}\right)$.
C'est bien la m\'ethode de H\'eron.\hfill$\square$

\textbf{(2--3)}
\begin{lstlisting}[language=Python]
def newton_sqrt(a, tol=1e-8):
    u = a / 2.0
    for _ in range(1000):
        u_new = 0.5 * (u + a / u)
        if abs(u_new - u) < tol:
            return u_new
        u = u_new
    return u

for a in [2, 3, 10]:
    print(f"sqrt({a}) ~ {newton_sqrt(a):.10f}")
\end{lstlisting}

\textbf{(4)} Newton converge quadratiquement (chaque it\'eration double les d\'ecimales exactes),
alors que la dichotomie converge lin\'eairement ($\log_2(1/\varepsilon)$ it\'erations).
En pratique~: Newton atteint $10^{-8}$ en $\sim5$ it\'erations,
la dichotomie en $\sim27$.

\needspace{3\baselineskip}
